\chapter*{Bibliographie et sources}
\addcontentsline{toc}{chapter}{Bibliographie et sources}
\markboth{Bibliographie et sources}{Bibliographie et sources}

\begin{small}
\section*{Statistiques et sources de données citées}
\begin{itemize}[leftmargin=1em]
\item Insee (2026), « Les comptes de la Nation en 2025 : le PIB ralentit et le pouvoir d'achat des ménages se replie »,
\emph{Insee Première} n\textsuperscript{o} 2105, mai 2026 ; et « Comptes nationaux annuels en 2025 », \emph{Insee
Résultats}, juin 2026. \url{https://www.insee.fr/fr/statistiques/8996855}
\item Insee (2026), « En 2025, le déficit public s'élève à 5,1~\% du PIB, la dette publique à 115,6~\% du PIB »,
\emph{Informations rapides} n\textsuperscript{o} 78, 27 mars 2026. \url{https://www.insee.fr/fr/statistiques/8956575}
\item Insee (2026), « Le compte des administrations publiques en 2025 », \emph{Insee Première} n\textsuperscript{o} 2106,
mai 2026. \url{https://www.insee.fr/fr/statistiques/8997691}
\item Insee (2026), « Niveau de vie et pauvreté en 2024 », \emph{Insee Première} n\textsuperscript{o} 2117, juillet 2026.
\url{https://insee.fr/fr/statistiques/9019316}
\item Insee (2026), « Taux de chômage au sens du BIT », \emph{Informations rapides}, 4\textsuperscript{e} trimestre 2025 et
1\textsuperscript{er} trimestre 2026 ; « Une photographie du marché du travail en 2025 », \emph{Insee Première}
n\textsuperscript{o} 2096. \url{https://www.insee.fr/fr/statistiques/8901327}
\item Insee (2026), « Comptes nationaux trimestriels, résultats détaillés du premier trimestre 2026 »,
\emph{Informations rapides} n\textsuperscript{o} 133.
\item Banque de France (2026), \emph{Rapport annuel 2025 de la balance des paiements et de la position extérieure de la
France}. \url{https://www.banque-france.fr/fr/communiques-de-presse/rapport-annuel-2025-de-la-balance-des-paiements-un-deficit-courant-quasi-stable-116-milliards-deuros}
\item Banque de France (2025-2026), balance des paiements mensuelle de la France (décembre 2025).
\item SDES (2026), \emph{Bilan énergétique de la France en 2025 — données provisoires}, avril 2026.
\url{https://www.statistiques.developpement-durable.gouv.fr/bilan-energetique-de-la-france-en-2025-donnees-provisoires}
\item FMI (2026), \emph{World Economic Outlook}, octobre 2026 (API DataMapper).
\item Conseil d'orientation des retraites (2026), \emph{Rapport annuel}, juin 2026.
\item Haut Conseil des finances publiques (2026), avis relatif au projet de loi d'approbation des comptes 2025, avril 2026.
\item Insee, projections de population 2026 ; recensement de la population 2022, fichier détail « Individus localisés au
canton-ou-ville ».
\item DGFiP, IRCOM 2025 (revenus 2024) ; statistiques nationales sur l'impôt sur le revenu par case.
\item Eurostat, bases de données de comptabilité nationale, d'entrées-sorties, de finances publiques et d'enquête sur
les forces de travail (codes en annexe~\ref{ann:repro}).
\end{itemize}

\section*{Littérature économique}
\begin{itemize}[leftmargin=1em]
\item Arrondel, L., Lamarche, P., Savignac, F. (2015), « Wealth effects on consumption across the wealth distribution:
empirical evidence », Banque de France, document de travail n\textsuperscript{o} 552.
\item Bach, L., Bozio, A., Guillouzouic, A., Malgouyres, C. (2023), « Dividend taxes and the allocation of capital: evidence
from the 2018 French reform », Institut des politiques publiques.
\item Bom, P., Ligthart, J. (2014), « What have we learned from three decades of research on the productivity of public
capital? », \emph{Journal of Economic Surveys}, 28(5).
\item Brülhart, M., Gruber, J., Krapf, M., Schmidheiny, K. (2022), « Behavioral responses to wealth taxes: evidence from
Switzerland », \emph{American Economic Journal: Economic Policy}, 14(4).
\item Cabannes, P.-Y., Houdré, C., Landais, C. (2014), « Comment le revenu imposable des ménages aisés réagit-il à sa
taxation ? », \emph{Économie et Statistique}, n\textsuperscript{o} 467-468.
\item Campolongo, F., Cariboni, J., Saltelli, A. (2007), « An effective screening design for sensitivity analysis of large
models », \emph{Environmental Modelling \& Software}, 22(10).
\item Engle, R., Granger, C. (1987), « Co-integration and error correction: representation, estimation and testing »,
\emph{Econometrica}, 55(2).
\item Fagereng, A., Holm, M., Natvik, G. (2021), « MPC heterogeneity and household balance sheets », \emph{American
Economic Journal: Macroeconomics}, 13(4).
\item Garbinti, B., Goupille-Lebret, J., Piketty, T. (2018), « Income inequality in France, 1900-2014 », \emph{Journal of
Public Economics}, 162.
\item Garbinti, B., Lamarche, P. (2014), « Qui épargne ? Qui désépargne ? », \emph{Économie et Statistique},
n\textsuperscript{o} 472-473.
\item Insee (2024), « Les inégalités de consommation et d'épargne », \emph{Insee Focus} n\textsuperscript{o} 338.
\item Insee, DG Trésor (2024), « Le modèle Mésange : réestimation et nouvelles variantes », document de travail.
\item Jakobsen, K., Jakobsen, K., Kleven, H., Zucman, G. (2020), « Wealth taxation and wealth accumulation: theory and
evidence from Denmark », \emph{Quarterly Journal of Economics}, 135(1).
\item Jappelli, T., Pistaferri, L. (2014), « Fiscal policy and MPC heterogeneity », \emph{American Economic Journal:
Macroeconomics}, 6(4).
\item Kleven, H., Landais, C., Saez, E. (2013), « Taxation and international migration of superstars: evidence from the
European football market », \emph{American Economic Review}, 103(5).
\item Kleven, H., Landais, C., Saez, E., Schultz, E. (2014), « Migration and wage effects of taxing top earners »,
\emph{Quarterly Journal of Economics}, 129(1).
\item Kramarz, F., Philippon, T. (2001), « The impact of differential payroll tax subsidies on minimum wage employment »,
\emph{Journal of Public Economics}, 82(1).
\item Lehmann, E., Marical, F., Rioux, L. (2013), « Labor income responds differently to income-tax and payroll-tax
reforms », \emph{Journal of Public Economics}, 99.
\item Lemoine, M., Turunen, H., et al. (2019), « The FR-BDF model and an assessment of monetary policy transmission in
France », Banque de France, document de travail n\textsuperscript{o} 736.
\item Lucas, R. (1976), « Econometric policy evaluation: a critique », \emph{Carnegie-Rochester Conference Series on
Public Policy}, 1.
\item Morris, M. (1991), « Factorial sampling plans for preliminary computational experiments », \emph{Technometrics},
33(2).
\item Saez, E. (2001), « Using elasticities to derive optimal income tax rates », \emph{Review of Economic Studies}, 68(1).
\item Seim, D. (2017), « Behavioral responses to wealth taxes: evidence from Sweden », \emph{American Economic Journal:
Economic Policy}, 9(4).
\end{itemize}

\section*{Documentation du projet}
\begin{itemize}[leftmargin=1em]
\item Documentation des modules : \code{doc/M1\_demographie.md}, \code{M3\_cadre\_comptable.md},
\code{M4\_population.md}, \code{M5\_productif.md}, \code{M7\_M8\_comportements\_et\_bouclage.md},
\code{M8\_reference.md}, \code{M9\_activite\_retraites.md}, \code{M10\_couplage\_micro\_macro.md},
\code{M11\_europe.md}, \code{M13\_comportements.md}.
\item Audit des mécanismes (v0.5) : \code{doc/12\_audit\_mecanismes.md} ; contre-expertise (v0.6) :
\code{doc/13\_contre\_expertise.md} ; registre des paramètres : \code{doc/14\_registre\_parametres.md}.
\item Validations : \code{outputs/validation\_M*.md}, \code{outputs/validation\_hors\_echantillon.md},
\code{outputs/validation\_additivite.md}.
\end{itemize}
\end{small}
